% Part of main.tex; compile main.tex.
% ---------------------------------------------------------------------------
\begin{abstract}
Asynchronous I/O is widely used to hide the cost of writing checkpoints and
analysis output in high-performance computing (HPC) applications. The
application issues its I/O early, and a background engine completes it while
the application continues computing. However, much of the evidence for this
benefit comes from I/O kernels and benchmarks that emulate the computation
phase with a call to \texttt{sleep()}. A sleeping thread does not use its
core, its memory bandwidth, or the network, so it removes the resource
contention that determines whether the I/O can actually be hidden. In this
paper, we survey 17 evaluations and benchmarks spanning HDF5, ADIOS, PDC,
POSIX-level checkpointing, and deep-learning I/O, and find that
\texttt{sleep}, an unspecified compute phase, or no compute phase at all is
the norm. We then present an open benchmark that drives POSIX, MPI-IO, HDF5,
ADIOS2, and PDC with interchangeable compute surrogates calibrated to the same
standalone duration, including CPU-bound, memory-bound, and
communication-bound kernels, a stencil mini-app, and GPU kernels. We find
that the contention does not always appear in the same place. Depending on
the I/O library, either the computation slows down while the background
engine runs, or the background engine makes little progress and the
application later waits for the remaining I/O. In both cases, the run time
the application observes is longer than a \texttt{sleep}-based evaluation
predicts. On a 48-core node writing 120\,GiB per run, \ExcLongerFive of \ExcN
configurations take at least 5\% longer than \texttt{sleep} predicts, and up
to \ExcMaxPct\% longer. In addition, the speedup of the HDF5 asynchronous VOL
connector drops from \SpHdfSleep$\times$ with \texttt{sleep} to
\SpHdfHalo$\times$ with a halo exchange, and non-blocking MPI-IO is slower
than blocking MPI-IO in every configuration we measured. Even PDC, whose
server runs on a reserved core, runs up to \ExcmaxpctPdc\% longer than
\texttt{sleep} predicts, because a polling progress thread in each client
shares the core of the rank. Lastly, we provide
recommendations for evaluating asynchronous I/O and for interpreting
existing results that rely on \texttt{sleep}.
\end{abstract}
